\renewcommand{\bibname}{参考文献}

\begin{thebibliography}{99}
  \addcontentsline{toc}{chapter}{参考文献}
  \bibitem{shimada2020}
  情報処理学会 出版委員会 監修，嶋田 義皓，
  『量子コンピューティング 基本アルゴリズムから量子機械学習まで』，
  オーム社，2020．

  \bibitem{tsukano2023}
  束野 仁政，
  『量子コンピュータの頭の中：計算しながら理解する量子アルゴリズムの世界』，
  技術評論社，2023．

  \bibitem{schuld2020}
  Maria Schuld，Francesco Petruccione 著，大関 真之 監訳，
  荒井 俊太，篠島 匠人，高橋 茶子，御手洗 光祐，山城 悠 訳，
  『量子コンピュータによる機械学習』，共立出版，2020．

  \bibitem{nakahara2016}
  中原 幹夫，
  『量子物理学のための線形代数 ベクトルから量子情報へ』，
  培風館，2016．

  \bibitem{risalc-normal}
  理数アラカルト，「正規行列の性質」（2026年11月23日閲覧）．\\
  \texttt{https://risalc.info/src/normal-matrix.html}

  \bibitem{takeuchi-normal}
  武内 修，『線形代数II』「正規行列の対角化可能性」，武内＠筑波大（2026年11月23日閲覧）．\\
  \texttt{\noautoxspacing\detokenize{https://dora.bk.tsukuba.ac.jp/~takeuchi/?線形代数}\kern0pt\detokenize{II/正規行列の対角化可能性}}
\end{thebibliography}
